% Source: https://artofproblemsolving.com/community/c3h1516458_help_me
% Stable identity: aops-topic-1516458-q01
Let the hexagon have side $s$ and centre $O$. The region Fido can reach is the disc of radius $r=OM$, where $M$ is the midpoint of a side; this disc is the inscribed circle of the hexagon, so it lies inside the yard.

\begin{asy}
usepackage("amsmath");
size(5cm);
defaultpen(fontsize(10pt));
pair[] H; for(int i=0;i<6;++i) H.push(dir(60*i));
fill(circle((0,0),sqrt(3)/2),rgb(0.8,0.9,1));
draw(H[0]--H[1]--H[2]--H[3]--H[4]--H[5]--cycle,linewidth(1));
draw(circle((0,0),sqrt(3)/2),blue+linewidth(0.9));
pair M=(H[0]+H[1])/2;
draw((0,0)--M,dashed);draw((0,0)--H[0],dashed+gray(0.5));
real s=0.07; pair a=unit((0,0)-M)*s, c=unit(H[0]-M)*s; draw(M+a--M+a+c--M+c);
dot((0,0));label("$O$",(0,0),W);
label("$r$",(0,0)+0.5*(M-(0,0)),NW);label("$s$",(0,0)+0.55*H[0],S);
label("$M$",M,NE);
label("\mbox{}",(min(currentpicture,user=true)+max(currentpicture,user=true))/2,black); // keeps the text after the figure black when LaTeX typesets the labels inline
\end{asy}

The hexagon consists of six equilateral triangles of side $s$, and $OM$ is the height of one of them: $r=\dfrac{\sqrt3}{2}s$. Hence
\[
\text{disc}=\pi r^2=\dfrac{3\pi s^2}{4},\qquad \text{hexagon}=6\cdot\dfrac{\sqrt3}{4}s^2=\dfrac{3\sqrt3}{2}s^2 ,
\]
and the fraction is
\[
\dfrac{\dfrac{3\pi}{4}}{\dfrac{3\sqrt3}{2}}=\dfrac{\pi}{2\sqrt3}=\dfrac{\sqrt3}{6}\pi .
\]
So $a=3$, $b=6$ and $ab=\boxed{18}$.
